\scanpage{33}
\mspage{7}
Nous allons maintenant examiner de plus près le cas où \deleted{$H^+$ \unsure{invarie}} $\operatorname{Herm}^{*+}(V)^{H^+}$ n'est pas réduit à une demi-droite. On vérifie aisément que cela revient à dire que $H^+$ \deleted{\unsure{est tel}} fixe deux pts distincts $z,z'$ de $X$, donc est contenu dans le sous-groupe de $G^+$ \deleted{\unsure{défini}} formé des \unsure{pts} qui fixent $z,z'$, lequel est isom. à $\mathbb C^*$, \deleted{\unsure{et est}}\inserted{et admet un} unique sous-groupe cpt maximal \deleted{\unsure{$T$ de dim. $1$}}\inserted{$T$ isom. à $\mathbb U$.}
\deleted{Dans ce cas les métriques invariantes par}

Si $H^+\ne\{1\}$,\inserted{[alors $\{z,z'\}$ \unsure{donc|\deleted{de}} $T$ est uniquement déterminé]} les métriques invariantes par $H^+$ sont aussi celles invariantes par $T\supset H^+$,
\margin{\uline{à vérifier}}
et elles sont en corr. $1$-$1$ \deleted{\unsure{avec}} les \deleted{trajectoires}\inserted{orbites} de $T$ dans $X^*=X\setminus\{z,z'\}$, [NB $T$ opère librement sur $X^*$\deleted{\unsure{$z$}}].

La relation s'obtient en associant à tte métrique invariante par $T$, le \deleted{\unsure{point}} grand cercle ayant comme pôles $z,z'$, lequel est une \deleted{trajectoire}\inserted{orbite} sous $T$. En particulier, tte orbite\inserted{de $X^*$} sous $H^+$ (étant contenue dans une unique orbite de $T$) définit une telle métrique.

Si $H=H^+$, il n'y a \deleted{pas} rien à ajouter. Dans le
% archive foliation: 35
